\section{Conclusion}
\label{sec:conclusion}

We have introduced Centroidal Voronoi Districts (\cvd{}), a Structure-layer
algorithm that assigns census tracts to district seeds by graph BFS distance
and iterates seeds to the population-weighted graph-distance medoid of their
assigned districts.
\cvd{} is the discrete redistricting analogue of Lloyd's algorithm for
optimal quantisation~\citep{lloyd1982,du1999}, adapted to weighted planar
graphs via the graph-distance medoid as the centroid surrogate.

\paragraph{Position in the B-series algorithm landscape.}
\cvd{} is a geometric structure algorithm, complementing the graph-topology
structure algorithms (\metis{}, \sa{}, \geosec{}, \areasec{}).
The two families serve different purposes:
\begin{itemize}
\item \textbf{Graph-topology methods} (\metis{}, \sa{}): minimise boundary
  length (edge cut), producing plans that follow roads, rivers, and county
  boundaries.
  Best for: publication-quality single-plan generation when compactness is
  measured by boundary minimisation.
\item \textbf{Geometric methods} (\cvd{}, \bfsgrowth{}): assign by proximity,
  producing plans that fill geographic space radially.
  Best for: geometric baseline plans, compactness diversity in ensemble
  analysis, and ``most natural partition'' legal arguments.
\end{itemize}

The empirical compactness ordering observed in Section~\ref{sec:comparison} is:
\[
  \text{METIS} \;<\; \text{CVD(iters=20)} \;<\; \text{SA(steps=10)},
\]
where $<$ denotes lower PP (worse compactness).
\cvd{} sits between \metis{} and \sa{} in compactness, at approximately
\metis{} runtime (1.6$\times$ slower on NC vs.\ 12$\times$ for \sa{}).

\paragraph{When to use CVD.}
\begin{itemize}
\item \textbf{Geometric baseline}: \cvd{} provides the ``most geographically
  natural partition'' --- the Voronoi tessellation of census tracts around
  their nearest district seed.
  This is a legally defensible baseline for redistricting challenges.
\item \textbf{Ensemble diversity}: Adding \cvd{} plans to a redistricting
  ensemble enriches the distribution with geometrically compact plans that
  are structurally different from edge-cut-minimising plans.
\item \textbf{Fast initialisation}: At $O(n_{\mathrm{iter}} \times m)$ runtime,
  \cvd{} is the fastest compactness-improving structure method in the B-series.
  For exploratory runs, \cvd{} provides a valid, balanced, compact plan
  at a cost close to \metis{}.
\end{itemize}

\paragraph{Phase~2 geographic CVD.}
Section~\ref{sec:phase2} previewed Phase~2 \cvd{} using Euclidean distance
on EPSG:5070 projected tract centroids.
Phase~2 is expected to yield a 5--10\% further PP improvement for coastal
and elongated states where BFS hop count diverges from geographic distance.
Implementation requires adding \texttt{tract\_centroids} to
\texttt{LoadedGraph}, which is a Phase~2 data pipeline task.

\paragraph{Implementation.}
\cvd{} bisection is implemented in the \texttt{redist} binary under
\texttt{--structure centroidal-voronoi}.
The primary tunable parameter is \texttt{--cvd-iters} (default 20).
The YAML configuration format:
\begin{verbatim}
algorithm:
  structure: centroidal-voronoi
  cvd_iters: 20
  cvd_metric: graph-distance
  balance_tolerance: 0.5
\end{verbatim}

\paragraph{Open questions.}
Several extensions are identified for future work:
\begin{enumerate}
\item \textbf{Phase~2 geographic CVD} (Section~\ref{sec:phase2}): implement
  Euclidean Voronoi once \texttt{tract\_centroids} is added to
  \texttt{LoadedGraph}.
  Validate on Florida and Louisiana for the largest expected improvements.

\item \textbf{CVD + BisectionEnsemble}: run \be{}$(T=50)$ with \cvd{}
  bisectors to combine geometric proximity with multi-seed selection.
  This is expected to improve PP over single-seed \cvd{} by selecting
  the most compact geometric Voronoi plan across 50 initialisations.

\item \textbf{CVD as SA warm start}: use the \cvd{} plan as the
  initialisation for \sa{} bisection instead of the \metis{} plan.
  Since \cvd{} provides a geometrically compact starting point, \sa{}
  refinement from \cvd{} may find lower-EC solutions with different
  compactness properties than \sa{} refinement from \metis{}.

\item \textbf{$k > 2$ CVD}: The current implementation bisects ($k=2$)
  at each node.
  A direct $k$-way \cvd{} that places $k$ seeds simultaneously is a natural
  extension, applicable at nodes where the bisection tree strategy calls for
  a $k$-way split (as in the \texttt{nway} structure or prime-factor nodes
  with $k_v = p$ for odd primes $p$).

\item \textbf{Convergence for approximate medoid}: Formalise the convergence
  analysis for the approximate medoid (probe size 50).
  Empirical evidence suggests monotone descent in practice; a formal bound
  on the approximation error's effect on convergence rate would complete
  the theoretical picture.
\end{enumerate}

These extensions inherit the clean three-layer compositor design and require
no structural changes to the codebase beyond the Phase~2 data field addition.
\cvd{} provides a principled geometric complement to the graph-topology
structure methods in the B-series, rounding out the algorithm portfolio
with a method that fills geographic space by proximity rather than by
boundary minimisation.
